% !TeX root = ../main.tex
%%% ---------------------------------------------------------------------------
%%% 摘要（technical report 用；实验报告一般不需要，main.tex 里默认注释掉了）
%%%
%%% GB/T 7713.3 对科技报告摘要的要求与论文一致：
%%%   报道性摘要，能独立于全文存在，不自指、不引文献编号。
%%%   结构：目的 → 方法 → 结果（带数据）→ 结论
%%% 关键词写在 setup.tex 的 \keywords{} 里，这里的环境会自动排出来。
%%% ---------------------------------------------------------------------------

\begin{abstract}
为比较常用优化器在图像分类任务上的收敛特性，在 ResNet-18 与 \dataset{}
上对 SGD（含动量）、Adam、AdamW 三种优化器各做 5 次独立重复训练，
按 A 类不确定度评定给出各项指标的均值与标准不确定度。结果表明，
测试准确率为 SGD-M $\qty{94.82(7)}{\percent}$、AdamW
$\qty{94.56(6)}{\percent}$、Adam $\qty{93.21(7)}{\percent}$，
差异均超出随机波动；而达到 \qty{90}{\percent} 准确率所需轮数则以
Adam 类方法为优（约 19 轮对 32 轮）。学习率扫描显示 SGD-M 对学习率的
敏感性显著高于 Adam 类方法。结论是自适应方法以约 \qty{1}{\percent} 的
准确率上限换取了收敛速度与超参数鲁棒性，适用于调参预算受限的场景。
\end{abstract}
